\documentclass[10pt,a4paper]{report}

%---------------------------------------------------------------------------
%	REQUIRED PACKAGES AND  CONFIGURATIONS
%---------------------------------------------------------------------------

% PACKAGES FOR TITLES
\usepackage{titlesec}
\usepackage{color}

% PACKAGES FOR LANGUAGE AND FONT
\usepackage[utf8]{inputenc}
\usepackage[english]{babel}
\usepackage[T1]{fontenc} % Font encoding
\usepackage{kpfonts}
\tolerance=1
\emergencystretch=\maxdimen
\hyphenpenalty=10000
\hbadness=10000
\vfuzz=30pt

% PACKAGES FOR IMAGES
\usepackage{graphicx}
\graphicspath{{Images/}} % Path for images' folder
\usepackage{eso-pic} % For the background picture on the title page
\usepackage{subfig} % Numbered and caption subfigures using \subfloat
\usepackage{caption} % Coloured captions
\usepackage{transparent}
\usepackage{wrapfig}

% STANDARD MATH PACKAGES
\usepackage{amsmath}
\usepackage{amsthm}
\usepackage{bm}
\usepackage[overload]{empheq}  % For braced-style systems of equations
\usepackage{matlab-prettifier}
\usepackage{gensymb}

% PACKAGES FOR TABLES
\usepackage{tabularx}
\usepackage{longtable} % tables that can span several pages
\usepackage{colortbl}
\usepackage{xltabular}
\usepackage{makecell}

% PACKAGES FOR ALGORITHMS (PSEUDO-CODE)
\usepackage{algorithm}
\usepackage{algpseudocode}
\newcommand{\algorithmautorefname}{Algorithm}

% PACKAGES FOR REFERENCES & BIBLIOGRAPHY
\usepackage{hyperref} % Adds clickable links at references
\hypersetup{
	colorlinks=true,
	linkcolor=bluePoli,
	anchorcolor=black,
	citecolor=bluePoli,
	filecolor=black,
	menucolor=black,
	runcolor=black,
	urlcolor=black,
}
\usepackage{cleveref}
\usepackage[super,square,numbers,sort&compress]{natbib}
\bibliographystyle{unsrt}
\usepackage{tocbibind}

% PACKAGES FOR THE APPENDIX
\usepackage[page]{appendix}
\renewcommand{\appendixpagename}{\Large\color{bluePoli}Appendice}

% PACKAGES FOR ITEMIZE & ENUMERATES 
\usepackage{enumitem}

% OTHER PACKAGES
\usepackage{amsthm,thmtools,xcolor} % Coloured "Theorem"
\usepackage{comment} % Comment part of code
\usepackage{fancyhdr} % Fancy headers and footers
\usepackage{lipsum} % Insert dummy text
\usepackage{tcolorbox} % Create coloured boxes (e.g. the one for the key-words)
\usepackage{stfloats} % Correct position of the tables
\usepackage{cuted}
\usepackage{multicol}

%-------------------------------------------------------------------------
%	NEW COMMANDS DEFINED
%-------------------------------------------------------------------------

\newcommand{\bea}{\begin{eqnarray}} % Shortcut for equation arrays
\newcommand{\eea}{\end{eqnarray}}
\newcommand{\e}[1]{\times 10^{#1}}  % Powers of 10 notation
\newcommand{\mathbbm}[1]{\text{\usefont{U}{bbm}{m}{n}#1}} % From mathbbm.sty
\newcommand{\pdev}[2]{\frac{\partial#1}{\partial#2}}

\DeclareMathOperator{\atantwo}{atan2}

\newcommand{\twofig}[5]{
\begin{figure}[H]
	\begin{minipage}{0.48\linewidth}
		\centering
		\includegraphics[width=#5\linewidth]{#1}
	\end{minipage}\hfill
	\begin{minipage}{0.48\linewidth}
		\centering
		\includegraphics[width=#5\linewidth]{#2}
	\end{minipage}
	\caption{#4}
	\label{fig:#3}
\end{figure}
}

\newcommand{\twofigII}[7]{
\begin{figure}[H]
	\begin{minipage}{0.48\linewidth}
		\centering
		\includegraphics[width=#7\linewidth]{#1}
		\caption{#2}
		\label{fig:#3}
	\end{minipage}\hfill
	\begin{minipage}{0.48\linewidth}
		\centering
		\includegraphics[width=#7\linewidth]{#4}
		\caption{#5}
		\label{fig:#6}
	\end{minipage}
\end{figure}
}

\newcommand{\IVfig}[7]{
\begin{figure}[H]
	\begin{minipage}{0.48\linewidth}
		\centering
		\includegraphics[width=#7\linewidth]{#1}
	\end{minipage}\hfill
	\begin{minipage}{0.48\linewidth}
		\centering
		\includegraphics[width=#7\linewidth]{#2}
	\end{minipage}
	\vfill
	\begin{minipage}{0.48\linewidth}
		\centering
		\includegraphics[width=#7\linewidth]{#3}
	\end{minipage}\hfill
	\begin{minipage}{0.48\linewidth}
		\centering
		\includegraphics[width=#7\linewidth]{#4}
	\end{minipage}
	\caption{#6}
	\label{fig:#5}
\end{figure}
}

\newcommand{\cfig}[4]{
\begin{figure}[H]
	\centering
	\includegraphics[width=#4\linewidth]{#1}
	\caption{#2}
	\label{fig:#3}
\end{figure}
}

\newcommand{\rfig}[4]{
\begin{wrapfigure}{r}{#4\linewidth}
	\centering
	\includegraphics[width=\linewidth]{#1}
	\caption{#2}
	\label{fig:#3}
	\vspace*{-13pt}
\end{wrapfigure}
}

\newcommand{\lfig}[4]{
\begin{wrapfigure}{l}{#4\linewidth}
	\centering
	\includegraphics[width=\linewidth]{#1}
	\caption{#2}
	\label{fig:#3}
	\vspace*{-13pt}
\end{wrapfigure}
}

% Second chapter commands
\newcommand{\reddashedline}{\textcolor{red}{---}}
\newcommand{\orangedashedline}{\textcolor{orange}{---}}
\newcommand{\yellowdashedline}{\textcolor{yellow}{---}}
\newcommand{\greendashedline}{\textcolor{green}{---}}
\newcommand{\bluedashedline}{\textcolor{blue}{---}}
\newcommand{\purpledashedline}{\textcolor{purple}{---}}
\newcommand{\cyandashedline}{\textcolor{cyan}{---}}
\newcommand{\magentadashedline}{\textcolor{magenta}{---}}

%-------------------------------------------------------------------------
%	MAIN BODY OF THE REPORT
%-------------------------------------------------------------------------

% Set the geometric layout of the document
\usepackage{geometry}
\geometry{
  top = 2.5cm,
  left = 1.5cm,
  right = 1.5cm,
  bottom = 2.5cm,
  headheight = 2.2cm,
  headsep = 0cm,
}
\raggedbottom

% Create color bluePoli (-> manuale grafica coordinata:  https://www.polimi.it/fileadmin/user_upload/il_Politecnico/grafica-coordinata/2015_05_11_46xy_manuale_grafica_coordinata.pdf)
\definecolor{bluePoli}{cmyk}{0.4,0.1,0,0.4}

% Custom theorem environments
\declaretheoremstyle[
  headfont=\color{bluePoli}\normalfont\bfseries,
  bodyfont=\color{black}\normalfont\itshape,
]{colored}

\captionsetup[figure]{labelfont={color=bluePoli}} % Set colour of the captions
\captionsetup[table]{labelfont={color=bluePoli}} % Set colour of the captions
%\captionsetup[algorithm]{labelfont={color=bluePoli}} % Set colour of the captions

\theoremstyle{colored}
\newtheorem{theorem}{Theorem}[section]
\newtheorem{proposition}{Proposition}[section]

% Enhances the features of the standard "table" and "tabular" environments.
\newcommand\T{\rule{0pt}{2.6ex}}
\newcommand\B{\rule[-1.2ex]{0pt}{0pt}}

% Algorithm description
\newcounter{algsubstate}
\renewcommand{\thealgsubstate}{\alph{algsubstate}}
\newenvironment{algsubstates}{
    \setcounter{algsubstate}{0}%
    \renewcommand{\STATE}{%
    \stepcounter{algsubstate}%
    \Statex {\small\thealgsubstate:}\space}
    }{}
    
% Custom theorem environment
\newcolumntype{L}[1]{>{\raggedright\let\newline\\\arraybackslash\hspace{0pt}}m{#1}}
\newcolumntype{C}[1]{>{\centering\let\newline\\\arraybackslash\hspace{0pt}}m{#1}}
\newcolumntype{R}[1]{>{\raggedleft\let\newline\\\arraybackslash\hspace{0pt}}m{#1}}

% Custom itemize environment
\setlist[itemize,1]{label=$\bullet$}
\setlist[itemize,2]{label=$\circ$}
\setlist[itemize,3]{label=$-$}
\setlist{nosep}

% Create command for background pic
\newcommand\BackgroundPic{% Adding background picture
	\put(230,358){
		\parbox[b][\paperheight]{\paperwidth}{%
			\vfill
			\centering
			\transparent{0.4}
			\includegraphics[width=0.5\paperwidth]{raggiera_polimi.eps}%
			\vfill
}}}

% Set indentation
\setlength\parindent{0pt}

% Set abstract
\renewenvironment{abstract}{\section*{\abstractname}}{}

% Custom title commands
\titleformat{\chapter}
{\color{bluePoli}\normalfont\Huge\bfseries}
{\color{bluePoli}\thechapter.}{1em}{}
\titlespacing*{\chapter}
{0pt}{-15mm}{1ex}

\titleformat{\section}
{\color{bluePoli}\normalfont\Large\bfseries}
{\color{bluePoli}\thesection.}{1em}{}
\titlespacing*{\section}
{0pt}{2ex}{1ex}

\titleformat{\subsection}
{\color{bluePoli}\normalfont\large\bfseries}
{\color{bluePoli}\thesubsection.}{1em}{}
\titlespacing*{\subsection}
{0pt}{2ex}{1ex}

% Custom headers and footers
\pagestyle{fancy}
\fancyhf{}
      
\fancyfoot{}
\fancyfoot[C]{\thepage} % page
\renewcommand{\headrulewidth}{0mm} % headrule width
\renewcommand{\footrulewidth}{0mm} % footrule width

\makeatletter
\patchcmd{\headrule}{\hrule}{\color{black}\hrule}{}{} % headrule
\patchcmd{\footrule}{\hrule}{\color{black}\hrule}{}{} % footrule
\makeatother

% -> Create the header
\chead[C]{
\centering
\begin{tcolorbox}[arc=0pt, boxrule=0pt, colback=bluePoli!60, width=\textwidth, colupper=white]
    \textbf{Interplanetary trajectory design} \hfill \textbf{\author}  
\end{tcolorbox}
}

\renewcommand{\title}{Interplanetary trajectory design and planetary perturbations analysis}

\newcommand{\authors}{
\begin{tabular}{rll}
    10723712	&	Marcello Pareschi	&	(BSc Aerospace Engineering - Politecnico di Milano)\\
    10836125	&	Daniele Paternoster	&	(BSc Aerospace Engineering - Politecnico di Milano)\\
    10711624	&	Alex Cristian Turcu	&	(BSc Aerospace Engineering - Politecnico di Milano)\\
    10884250	&	Tamim Harun	Or		&	(BSc Aerospace Engineering - International Islamic University Malaysia)
\end{tabular}
}

\newcommand{\course}{Orbital Mechanics}

\newcommand{\professor}{Camilla Colombo}

%\newcommand{\firstcoadvisor}{Name Surname}
%\newcommand{\secondcoadvisor}{Name Surname}

\newcommand{\YEAR}{2023-2024}

\begin{document}
\pagenumbering{Roman}

%---------------------------------------------------------------------------
% TITLE PAGE
%---------------------------------------------------------------------------

% DO NOT REMOVE SPACES BETWEEN LINES!

\let\temp\newpage
\let\newpage\relax
\begin{titlepage}

\AddToShipoutPicture*{\BackgroundPic}

\hspace{-0.6cm}\includegraphics[width=0.6\textwidth]{logo_polimi_ing_indinf.eps}

\vspace{-0.2cm}
\Large{\textbf{\color{bluePoli}{\title}}}\\
\hspace*{\fill}

\vspace{-0.2cm}
\fontsize{0.3cm}{0.5cm}\selectfont \bfseries \textsc{\color{bluePoli} MSc in \course}\\
\hspace*{\fill}

\vspace{-0.2cm}
\fontsize{0.3cm}{0.5cm} \selectfont \bfseries Authors:\par
\vspace{0.2cm} \hspace{0.3cm} \textsc{\textbf{\authors}}\\
\hspace*{\fill}

\fontsize{0.3cm}{0.5cm}\selectfont \bfseries Professor: \textsc{\textbf{\professor}}\\
\hspace*{\fill}

% if only ONE co-advisor is present:
%\vspace{-0.4cm}
%\fontsize{0.3cm}{0.5cm}\selectfont \bfseries Co-advisor: %\textsc{\textbf{\firstcoadvisor}}\\
% if more than one co-advisors are present:
%\vspace{-0.4cm}
%\fontsize{0.3cm}{0.5cm}\selectfont \bfseries Co-advisors: \textsc{\textbf{\firstcoadvisor}}\textsc{\textbf{\secondcoadvisor}}\\

\vspace{-0.4cm}
\fontsize{0.3cm}{0.5cm}\selectfont \bfseries Academic year: \textsc{\textbf{\YEAR}}

\small \normalfont

\vspace{11pt}

\centerline{\rule{1.0\textwidth}{0.4pt}}

\vspace{15pt}

\end{titlepage}
\let\newpage\temp

\thispagestyle{plain} % In order to not show the header in the first page

\vspace*{20mm}
\begin{abstract}
\addcontentsline{toc}{chapter}{Abstract}
\vspace*{10mm}

\lipsum[1]

\end{abstract}

\clearpage

%---------------------------------------------------------------------------
% INDEX
%---------------------------------------------------------------------------

{\hypersetup{linkcolor=black}\tableofcontents}
\clearpage

%{\hypersetup{linkcolor=black}\listoftables}

%{\hypersetup{linkcolor=black}\listoffigures}

%---------------------------------------------------------------------------
% BEGINNING OF THE WRITTEN PART
%---------------------------------------------------------------------------

\setcounter{page}{1}
\pagenumbering{arabic}

\chapter{Interplanetary mission}
\section{Symbols}
\label{sec:symbols}
%\setcounter{subsection}{2}
\newsavebox\ltmcbox
\renewcommand{\arraystretch}{1.2}
\renewcommand\tabularxcolumn[1]{m{#1}}

\begin{multicols}{2}
	\raggedcolumns

	\subsubsection{Analisi della missione}
	\vspace{-10pt}
	\small
	\setbox\ltmcbox\vbox{
	\makeatletter\col@number\@ne
	\begin{xltabular}{\linewidth}{cc>{\raggedright\arraybackslash}X}
		$\bm{A_e}$ & $[m^2]$ & area di efflusso totale \\
		$\bm{\phi}$ & $[rad]$ & angolo di traiettoria del razzo
	\end{xltabular}
	\unskip
	\unpenalty
	\unpenalty}
	\unvbox\ltmcbox

	\subsubsection{Analisi della missione 2}
	\vspace{-10pt}
	\small
	\setbox\ltmcbox\vbox{
	\makeatletter\col@number\@ne
	\begin{xltabular}{\linewidth}{cc>{\raggedright\arraybackslash}X}
		$\bm{A_e}$ & $[m^2]$ & area di efflusso totale \\
		$\bm{\phi}$ & $[rad]$ & angolo di traiettoria del razzo
	\end{xltabular}
	\unskip
	\unpenalty
	\unpenalty}
	\unvbox\ltmcbox

\end{multicols}
\subsection{Introduction}
\label{subsec:introduction}

\subsubsection{Description of the problem}
\label{subsubsec:description}
The first part of the assignment aims at designing an interplanetary transfer from Mars to asteroid 1036 Ganymede exploiting a powered gravity assist on Earth. The problem is analyzed through the patched conics method, without considering the injection and arrival hyperbolas. The initial and final velocity vector of the satellite are assumed to be the same of the respective celestial body. The two heliocentric legs are calculated through the Lambert problem. The trajectory is selected with the only criteria of minimizing the cost of the mission, assessed through the total $\Delta V$. The latter is computed by summing different contributions:

\begin{equation}
    \Delta V_{tot}= \Delta V_1 + \Delta V_2 + \Delta V_3
\end{equation}

Where the three terms are defined as:
\begin{itemize}
    [wide,itemsep=3pt,topsep=3pt]
    \item $\Delta V_1$ injection in the first heliocentric leg;
    \item $\Delta V_2$ exit from the second heliocentric leg;
    \item $\Delta V_3$ impulse given by the engine at pericentre of hyperbola fly-by;
\end{itemize}


All the manoeuvres are assumed to be impulsive, i.e. they change only the velocity vector of the spacecraft, mantaining invariated the position vector. Note that $\Delta V_{1}$ and $\Delta V_{2}$ are related to heliocentric velocities, while $\Delta V_{3}$ is calculated through relative geocentric velocities. 

For each manoeuvre the cost is computed as:

\begin{equation}
    \Delta V_i= \lVert \boldsymbol{V_{+,i}} - \boldsymbol{V_{-,i}} \rVert
\end{equation}


\subsubsection{Assigned data and constraints}
\label{subsubsec:data_constraints}
A few constraints were considered:
\begin{itemize}
    [wide,itemsep=3pt,topsep=3pt]
    \item earliest departure date $t_{min,dep} = \left[01/01/2028 \;\; 00:00:00\right]$ in Gregorian calendar
    \item latest arrival date $t_{max,arr} = \left[01/01/2058 \;\; 00:00:00\right]$ in Gregorian calendar
    \item time of flights of the two heliocentric arcs must be greater than the associated parabolic time
    \item minimum pericentre radius of the fly-by hyperbola $r_p = r_E + 500 \; km$
    \item single-revolution Lambert problem was considered
    \item reasonable total time of the mission
\end{itemize}

Note that the constraint on the pericentre radius of the fly-by is considered both for avoid impact on Earth and to prevent undesired atmoshperic drag effects. 




\section{Algorithms description}
\label{sec:algo_description}

The targeting problem previously defined can be seen as an optimization problem with three degrees of freedom (DOFs). Indeed, once the departing date and the two times of flight of the heliocentric legs are chosen, both the Lambert's arcs and the fly-by hyperbola are fully defined.  
Regarding the formulation of the Lambert's problem, it requires the knowledge of the initial and final position and also the imposed time of flight between them. For two Lambert's arcs it would be needed a total of six DOFs, but this quantities are dependant one to each other:
\begin{itemize}
    [wide,itemsep=3pt,topsep=3pt]
    \item the final position vector of the first arc corresponds to the initial position of the second one;
    \item the initial date for the departing on the second arc corresponds the arrival date on the first arc;
    \item fixing the first Lambert's arc, the final position and arrival date for the second Lambert's arc are related through the analytical ephemerides.
\end{itemize}

Once the two heliocentric legs are determined, the powered gravity assist follows as the geocentric velocity vectors are known.

The method implemented to solve the optimization problem is the \textbf{brute force algorithm} refined with the \textbf{gradient descent algorithm}. Then, to validate the results, the \textbf{brute force algorithm} has been used solely with a more dense search grid in a reasonable domain. 

\subsection{Refined brute force algorithm}
\label{subsec:brute_force_algo}

The presented \autoref{alg:brute} is reliable, yet  computationally demanding since the research of the minimum is performed through a triple-nested \textit{for} loop. The time of the execution highly depends on the refinement chosen for the three selected periods of time. From this considerations and also noticing that the function to minimize presents high irregularities, it is clear that narrow time windows can help in the overall research. Moreover, with this last observation, it is possible to use a less fine grid of research, since it is reasonable to think that the function will present less irregularities.

\begin{algorithm}
    \caption{Bruteforce algorithm} \label{alg:brute}
    \begin{algorithmic}
    \Require $T_{dep}, \Delta T_{1}, \Delta T_{2}$
    \State $\Delta V_{min} = 10^{10}$
    \For{$i$ in $T_{dep}$}
        \For{$j$ in $\Delta T_{1}$} 
            \State Calculate $T_{fly-by}$, first Lambert's arc, Mars' velocity at $T_{dep}$ and  $\Delta V_1$
            \For{$k$ in $\Delta T_{2}$}
                \State Calculate $T_{arr}$, second Lambert's arc, Asteroid's velocity at $T_{arr}$, $\Delta V_2$, $\Delta V_3$, $\Delta V_{tot}$
                \If{$\Delta v_{tot} < \Delta v_{min}$ and $r_p > r_{Earth} + 500 \; km$}
                \State $\Delta v_{min} = \Delta v_{tot}$
                \State $T_{dep,min} = T_{dep} $;   $T_{fb,min}  = T_{fb}$;   $T_{arr,min} = T_{arr}$
                \EndIf
            \EndFor
        \EndFor
    \EndFor
\State Minimize cost using \textit{fminunc} with initial guess  $\left(T_{dep,min};T_{fb,min};T_{arr,min} \right)$
\\
\Return $\Delta V_{min}$; $T_{dep,min}$; $T_{fb,min}$; $T_{arr,min}$
    \end{algorithmic}
\end{algorithm}

As a consequence, it was decided to perform a physical analysis (\autoref{sec:time_red}) in order to wisely reduce the time domains for departure, fly-by and arrival. The idea of \autoref{alg:brute} is that, once this reduction is performed, the brute-force search can be carried out on a small but refined grid. To speed up the process and refine the solution, \textit{fminunc} of \textit{MATLAB} was used. The selected initial guess is the outcome of the brute force research.

In order to completely validate the results, the robustness of the brute force algorithm has been exploited. Referring to \autoref{alg:brute}, \textit{fminunc} was removed after the for-loop search and the research grid was refined. Moreover, with this computation it is possible to validate the time reduction analyzed in \autoref{sec:time_red}.
\subsection{Reduction of the time windows}
\label{subsec:time_red}

\subsubsection{Resonance period analysis}
\label{subsubsec:res_period}

\subsubsection{Cost-plot analysis}
\label{subsubsec:cost_plot_analysis}

\subsubsection{Final time window selection}
\label{subsubsec:final_window}
\section{Conclusion and results}
\label{sec:conclusion}

It was then used \autoref{alg:brute} in the time domain selected at the end of \autoref{subsec:cost_plot_analysis},  using 40 elements for each array. The obtained results in terms of cost and dates are reported in the following tables:

\vspace{0.5cm}

\begin{minipage}{0.5\linewidth}
    \centering
    \captionsetup{type=table}
    \begin{tabular}{|c|c|c|c|}
        \hline
        $\Delta V_1 \, [km/s]$ & $\Delta V_2 \, [km/s]$ & $\Delta V_3 \, [km/s]$ & $\Delta V_{tot} \, [km/s]$\\
        \hline
        $4.0813$ & $3.3130$ & $2.7715 \cdot 10^{-6}$ & $7.3942$\\
        \hline
    \end{tabular}
    \caption{Costs $\Delta V$ of the optimal mission}
    \label{table:cost_optimal}
\end{minipage}\hfill
\begin{minipage}{0.5\linewidth}
    \centering
    \captionsetup{type=table}
    \begin{tabular}{|c|c|c|}
        \hline
        Dep. Date & Fb. Date & Arr. Date \\
        \hline
        $\left[03/06/2033 \right]$ & $\left[10/02/2034\right]$ & $\left[27/03/2036\right]$ \\
        \hline
    \end{tabular}
    \caption{Optimal mission dates in Gregorian calendar}
    \label{table:date_optimal}
\end{minipage}
\vspace*{5pt}

The obtained results were validated through the pure brute-force algorithm, with a more refined grid of search in the same domain. In particular 100 elements where considered, the data obtained were:

... put the two tables of brute force algorithm

\subsection{Heliocentric legs}
\label{subsec:heliocentric}
The two heliocentric legs associated with the solution computed through \autoref{alg:brute} can be characterized through the keplerian elements shown in \autoref{table:Orb_Par}:
\begin{table}[H]

    \centering
    \begin{tabular}{|c|c|c|c|c|c|c|c|}
    \hline
    Arc &  $a  \ [AU]$ & $e \ [-]$ & $i \ [deg]$ & $\Omega \ [deg]$ & $\omega \ [deg]$ & $\theta_{dep} \ [deg]$ & $\theta_{arr} \ [deg]$ \\
    \hline
    M$\to$E &  $1.1602$ & $0.2870$ & $1.18$ & $321.55$ & $105.97$ & $195.38$ & $74.03$ \\
    \hline
    E$\to$G & $1.9445$ & $0.5029$ & $4.17$ & $141.55$ & $339.91$ & $20.09$ & $231.37$ \\
    \hline
    \end{tabular}
    
    \caption{Heliocentric legs}
    \label{table:Orb_Par} 
\end{table}

\cfig{solar_system.jpg}{Optimal transfer trajectory - planet's dimension not in scale}{solar_system}{0.75}

The two heliocentric legs are displayed in \autoref{fig:solar_system}. The represented positions of Mars, Earth and 1036 Ganymed are respectively at departure, fly-by and arrival. All the orbits are plotted in the Heliocentric ecliptic inertial frame.
\subsection{Gravity assist}
\label{subsec:ga}

Known the incoming and outcoming heliocentric velocities of the spacecraft and the Earth's velocity at the fly-by date, it was possible to completely characterize the two branches of hyperbola in an Earth-centred ecliptic frame. The main parameters of the hyperbola are reported in \autoref{table:ga}:

\begin{table}[H]

    \centering
    \begin{tabular}{|c|c|c|c|c|c|c|c|}
    \hline
    $r_p \ [km]$ &  $\delta _{tot} \ [\deg]$ & $v_{\infty}^{-} \ [km/s]$ & $v_{\infty}^{+} \ [km/s]$ & $e_{-} \ [-]$ & $e_{+} \ [-]$ & $\Delta T_{soi} \ [days]$ & $\Delta V_{helio} \ [km/s]$ \\
    \hline
    $7763.39$ &  $54.8918$ & $7.743439$ & $7.743435$ & $2.167834$ & $2.167833$ & $2.6871$ & $7.1381$ \\
    \hline
    \end{tabular}
    
    \caption{Gravity assist data}
    \label{table:ga} 
\end{table}

It can be noticed that the fly-by hyperbola respect the constraint on the pericentre radius given in \autoref{subsec:data_constraints}. Moreover, the impulse given by the engine at the pericentre (\autoref{table:cost_optimal}) is negligible with respect to the heliocentric fly-by manoeuvre cost, calculated as $\Delta V_{helio} = \lVert\boldsymbol{V_{\infty}^{+} - V_{\infty}^{-}} \rVert$. This confirms the efficiency effect of the examined fly-by. 

The permanence time passed inside the sphere of influence has been computed using the hyperbolic time law. In particular, considering the intersection between the two hyperbolic arcs and the radius of the Earth's sphere of influence (\autoref{eq:soi}), the true anomalies at entry and exit of SOI are computed.

\begin{equation}
    \label{eq:soi}
r_{SOI,E} = r_{S \to E}\left(\frac{M_E}{M_S}\right)^{2/5} = 9.24647 \cdot 10^{5} \ km
\end{equation}













%---------------------------------------------------------------------------

\chapter{Planetary mission}


%---------------------------------------------------------------------------
% BIBLIOGRAFIA
%---------------------------------------------------------------------------

\clearpage
\bibliography{bibliography}

\end{document}